\section{Finiteness of the compact coefficient cohomology}
\label{sec:coefficients}

We now close the induction on the height strata.  Throughout this section
$p$ and $n$ are fixed, $G=P_n$, and $\kk$ is a finite field containing the
fields of definition used in the finite marking constructions.  Write
\[
 A_m=\kk[[u_m,\ldots,u_{n-1}]],\qquad 1\leq m\leq n,
\]
so that $A_n=\kk$.  The coefficient class is that of
\Cref{def:mark-coefficients}: integer powers of the periodicity line,
whole finite division-tuple algebras with the permitted Frobenius powers,
and their finite tensor products, together with the stated closure under
finite sums, summands, and finite filtrations.  All coefficients over
$A_m$ carry their compact topology.

\begin{theorem}\label{thm:coeff-finite}
For every $1\leq m\leq n$, every coefficient module $M$ in
\Cref{def:mark-coefficients} over $A_m$, and every $i\geq0$, the group
\[
 H^i_{\cts}(P_n,M)
\]
is finite.  In particular, this holds for every integer periodicity twist
of $A_m$.
\end{theorem}

The proof uses countability on the open stratum and compactness on the
closed formal neighborhood.  We first supply the coefficient and descent
steps that connect these two statements.  For $m<n$, put
$x=u_m$, $B=B_m=A_m[1/x]$, and $r=n-m$.  We choose finite marking levels
from a cofinal sequence of normal product subgroups
\[
 U\triangleleft J=P_m\times\GL_r(\Zp).
\]
Thus $B_U/B$ is a finite \'etale $J/U$-torsor.  Normality is available by
taking sufficiently deep congruence subgroups; it imposes no extra
hypothesis on the fixed group $G$.

\begin{lemma}\label[lemma]{lem:coeff-finite-representations}
Assume the coefficient finiteness assertion at height $m+1$.
For a basic coefficient $M$ over $A_m$, including any fixed tensor product
of the basic division-tuple coefficients and an integer periodicity
twist, there is a sufficiently deep normal product level $U$ such that
\[
 H^i_{\cts}\bigl(G,B_U\otimes_{A_m}M\bigr)
\]
is countable for every $i\geq0$.  The same conclusion holds after further
shrinking within the chosen cofinal sequence.
\end{lemma}

\begin{proof}
By \Cref{thm:pro-countability}, whose boundary hypothesis is precisely
\Cref{lem:mark-strips} at height $m+1$, all groups
$H^i_{\cts}(G,B_U)$ are countable on sufficiently deep fixed levels.
We show that the indicated coefficient extensions have the same property.

Apply \Cref{lem:mark-associated-coefficients} to the basic coefficient
$M$. Its whole division-tuple factors, including their powered
coordinates, give a finite-dimensional translation representation $E$;
the linear connected marking trivializes the chosen integer periodicity
power. At every sufficiently deep level this gives a $G$-equivariant
identification
\[
 B_U\otimes_{A_m}M\simeq V_E.
\]
Here $G$ acts on the generator-lift torsor and commutes with translation.

We will use a finite filtration or a finite direct-summand diagram of
$E$. Its finitely many coaction matrices and arrows factor through one
finite translation quotient. By the same finite-presentation descent
as in \Cref{lem:mark-associated-coefficients,sec:transfers}, the entire
diagram and its torsor descend to one further finite marking level.
The auxiliary actions on its finite-dimensional $\kk$-spaces factor
through finite groups, so a common shrink makes them trivial on the
diagram. Naturality under formal-group isomorphisms retains
$G$-equivariance: the identities hold at full markings and descend
faithfully to the fixed finite stage. These choices concern finitely
many modules and maps, not individual elements of $G$.

The associated-bundle functor is exact on these finite diagrams.
After base change to the finite lift torsor it is tensoring with $E$;
faithful flatness detects exactness.  Its terms are finite locally free
bundles on the affine open.  Surjections between them have underlying
module sections, so after choosing lattices their sections have bounded
denominators.  Consequently the associated exact sequences give exact
continuous cochain sequences in the convention of
\Cref{lem:cts-exact}.

Suppose first that $m>1$.  The distribution algebra is
\[
 \Lambda=\kk[[D_1,\ldots,D_d]],\qquad d=(m-1)r,
\]
with maximal ideal $\mathfrak m=(D_1,\ldots,D_d)$, as in
\Cref{prop:trans-roots}.  A finite translation representation is a
finite-length continuous $\Lambda$-module.  Its finite filtration
\[
 E\supseteq\mathfrak mE\supseteq\mathfrak m^2E
 \supseteq\cdots\supseteq0
\]
has quotients that are finite direct sums of the trivial representation
$\kk$.  Having retained this finite diagram at the chosen marking level,
its associated filtration has quotients $B_U^{\oplus a}$ with their
ordinary $G$-actions.  Countability of their cohomology and the long exact
sequences of the filtration give countability of
$H^i_{\cts}(G,V_E)$ in every degree.

For $m=1$, the translation group is of multiplicative type.  Choose
$H$ from the full cofinal sequence of finite generator-lift quotients,
deep enough that $E$ factors through $H$.  Thus its torsor is the full
root algebra of \Cref{prop:mark-roots}, rather than an intermediate
subalgebra of that root extension.  Its
distribution algebra
\[
 R_H=\mathcal O(H)^\vee=\mathcal O(H^D)
\]
is finite \'etale over $\kk$, because the Cartier dual $H^D$ of a finite
multiplicative-type group is \'etale.  Every finite $R_H$-module is
therefore a direct summand of $R_H^{\oplus a}$ for some $a$.
The perfect trace pairing of the finite \'etale algebra identifies
$R_H$ with $R_H^\vee=\mathcal O(H)$ as $R_H$-modules.  Thus $E$ is a
direct summand of a finite sum of regular function representations.
This argument works over $\kk$ itself and includes quotients whose
order is divisible by $p$.
For the general semisimplicity statement over the ground field, see
\citet[Theorem~12.30, pp.~241--242]{milne2017}.
Shrink $U$ so that the splitting maps of this finite diagram descend as well.

The associated bundle of such a regular representation is the function
algebra of the corresponding finite lift torsor.  By
\Cref{prop:mark-roots} and \eqref{eq:trans-full-torsor}, after enlarging
the level it is
\[
 B_U^{1/h},\qquad h=p^a.
\]
The exponent may be chosen to fix $\kk$.  Powering gives an additive,
$G$-equivariant homeomorphism
\begin{equation}\label{eq:coeff-root-powering}
 B_U^{1/h}\longrightarrow B_U,\qquad z\longmapsto z^h.
\end{equation}
For the bounded-pole topology, this follows either from the finite root
coordinates or by multiplying every root pole bound and parameter order
by $h$; the inverse rescales those bounds.  Hence powering identifies
the continuous cochain complexes as additive complexes.  It is also
$\kk$-linear for the chosen exponents.  It follows that each regular
block, and then its finite sums and summands, has countable cohomology.

All choices involve finitely many levels and finite diagrams.  They persist on a
cofinal tail, where \Cref{thm:pro-countability} also applies.  This proves
the assertion in both cases.
\end{proof}

\begin{lemma}\label[lemma]{lem:coeff-finite-descent}
Let $M_B$ be a finite locally free $B$-module with the coefficient action
of $G$, and let $U\triangleleft J$ be a finite marking level.  If
\[
 H^i_{\cts}(G,B_U\otimes_B M_B)
\]
is countable for every $i\geq0$, then $H^i_{\cts}(G,M_B)$ is countable for
every $i\geq0$.  The order of $J/U$ need not be invertible in $\kk$.
\end{lemma}

\begin{proof}
Put $H=J/U$ and $N=B_U\otimes_B M_B$.  The $G$- and $H$-actions on $N$
commute.  Finite \'etale torsor descent identifies the augmented Amitsur
complex with the finite-group cochain resolution
\begin{equation}\label{eq:coeff-marking-amitsur}
 M_B\longrightarrow N\longrightarrow C^1(H,N)
 \longrightarrow C^2(H,N)\longrightarrow\cdots,
\end{equation}
where the displayed $N$ is in group-cochain degree zero.  This augmented
complex is exact.  For completeness, after faithfully flat base change
to $B_U$ the torsor becomes the trivial $H$-torsor and its augmented
complex has the usual extra-degeneracy contraction.  Exactness therefore
holds before base change as well.

It also holds with the required continuous $G$-cochain parameters.
To make the topology explicit, choose
$\tau\in B_U$ with $\sum_{h\in H}h(\tau)=1$.  Such an element exists:
the torsor trace is surjective after the same faithfully flat base
change, so its cokernel vanishes over $B$.  The $B$-linear map
\[
 \lambda:B_U\longrightarrow B,\qquad
 \lambda(z)=\Tr_{B_U/B}(\tau z)
\]
satisfies $\lambda(1)=1$.  Applying $\lambda$ to the first algebra
factor contracts the augmented Amitsur complex; the insert-unit
differential gives the contraction identity directly.  These are
finite matrices over $B$ on the fixed finite locally free terms, so
they send each lattice to a bounded pole shift and are continuous
there.  The contraction need not commute with $G$: exactness of each
row after applying continuous cochains is what is needed.  Thus applying $C^a_{\cts}(G,-)$ to
\eqref{eq:coeff-marking-amitsur} remains exact for every $a$.

The first-quadrant double complex
\[
 C^a_{\cts}\bigl(G,C^b(H,N)\bigr)
\]
therefore computes $H^*_{\cts}(G,M_B)$, and its other spectral sequence is
\begin{equation}\label{eq:coeff-marking-spectral}
 H^b\bigl(H,H^a_{\cts}(G,N)\bigr)
 \ \Longrightarrow\ H^{a+b}_{\cts}(G,M_B).
\end{equation}
Finite $H$-cochains are finite products.  Their cohomology on a countable
module is countable, regardless of divisibility of $|H|$ by $p$.
Every total-degree diagonal of \eqref{eq:coeff-marking-spectral} is
finite, so its abutment is countable.  This proves the claim.
\end{proof}

\begin{proof}[Proof of \Cref{thm:coeff-finite}]
Tensoring finite free modules preserves exact filtrations and retracts,
and a tensor product of basic coefficients is basic. Induction on the
finite constructions in \Cref{def:mark-coefficients} therefore reduces
the asserted closure properties to finite filtrations and retracts.
We induct downwards on $m$.  At $m=n$, every basic coefficient is a
finite-dimensional vector space over the finite field $\kk$.  A
finite-type projective resolution for $G$ computes its cohomology by a
complex whose term in each degree is a summand of a finite power of this
finite coefficient module; see \Cref{lem:cts-comparison}.  Each
cohomology group is therefore finite.  This uses finiteness of the
resolution in each degree, not finite cohomological dimension of $G$.
Finite sums, summands, and finite filtrations preserve the conclusion.

Now let $m<n$ and assume the assertion at $m+1$.  It suffices first to
treat a basic tensor coefficient $M$: the long exact cohomology sequence
and retracts then give all the closures allowed in
\Cref{def:mark-coefficients}.  By
\Cref{lem:coeff-finite-representations,lem:coeff-finite-descent},
\begin{equation}\label{eq:coeff-open-countable}
 H^i_{\cts}(G,M[1/x])\quad\text{is countable for every }i\geq0.
\end{equation}

We spell out the boundary comparison.  The submodules $x^aM$ are
$G$-stable as submodules with their induced actions, since $G$ preserves
the ideal $(x)$.  The successive quotients of every fixed pole strip
are restrictions of the original coefficients to $A_{m+1}$, with the
appropriate conormal twists.  These are still allowed coefficients:
the conormal line is an integer periodicity twist, and a powered tuple
with $b\leq m\ell$ also satisfies $b\leq(m+1)\ell$.  The whole-model
filtrations in \Cref{prop:mark-whole-model,lem:mark-strips} consequently
give finite cohomology for every fixed strip.  In particular, for
\[
 Q_a=x^{-a}M/M\quad(a\geq0),\qquad
 Q=M[1/x]/M=\colim_a Q_a,
\]
each $H^i_{\cts}(G,Q_a)$ is finite.

By the bounded-denominator convention, cochains on $Q$ are precisely
the filtered union of the cochains on the $Q_a$.  Filtered colimits of
vector spaces are exact, so
\[
 H^i_{\cts}(G,Q)=\colim_a H^i_{\cts}(G,Q_a)
\]
is countable.  The exact cochain sequence supplied by
\Cref{lem:cts-exact} for $0\to M\to M[1/x]\to Q\to0$ now shows that the
kernel and cokernel of
\[
 H^i_{\cts}(G,M)\longrightarrow H^i_{\cts}(G,M[1/x])
\]
are countable.  Together with \eqref{eq:coeff-open-countable}, this proves
that $H^i_{\cts}(G,M)$ is countable.

Finally, $M$ is compact.  The finite-type completed resolution in
\Cref{lem:cts-comparison} computes $H^i_{\cts}(G,M)$ as a quotient of a
closed subspace of a compact module by the compact image of a continuous
map.  Its cohomology is thus compact and Hausdorff.  A countable compact
Hausdorff group is finite: by the Baire theorem one of its singleton
closed subsets has interior; translations make the group discrete, and
compactness makes that discrete group finite.  Hence
$H^i_{\cts}(G,M)$ is finite.  This closes the induction.
\end{proof}

\begin{corollary}\label[corollary]{cor:coeff-morava}
Let $E_n$ be Morava $E$-theory for the height-$n$ Honda group over
$\mathbb F_{p^n}$, and let
\[
 \mathbb G_n=P_n\rtimes\Gal(\mathbb F_{p^n}/\mathbb F_p)
\]
be its extended stabilizer.  For every $i\geq0$ and every $t\in\mathbb Z$,
both groups
\[
 H^i_{\cts}\bigl(P_n,\pi_t(E_n/p)\bigr),\qquad
 H^i_{\cts}\bigl(\mathbb G_n,\pi_t(E_n/p)\bigr)
\]
are finite.
\end{corollary}

\begin{proof}
The coefficient ring is
\[
 \pi_*E_n=W(\mathbb F_{p^n})[[u_1,\ldots,u_{n-1}]][u^{\pm1}],
 \qquad |u|=-2.
\]
Multiplication by $p$ is injective, and the ring is even.  Therefore the
odd homotopy groups of $E_n/p$ vanish, while each even homotopy group is
the special-fiber ring with an integer power of its intrinsic periodicity
line.  Extending $\mathbb F_{p^n}$ to the finite field $\kk$ identifies
this coefficient with one of the $m=1$ twists in
\Cref{thm:coeff-finite}.  Finite scalar extension commutes with the compact
continuous-cochain complex and with its cohomology: after choosing a
field basis it is a finite direct sum, with the product compact topology.
Faithfulness of the extension therefore proves the asserted finiteness
for $P_n$ over $\mathbb F_{p^n}$.

For the extended group use the continuous group-extension spectral
sequence
\[
 H^a\!\left(\Gal(\mathbb F_{p^n}/\mathbb F_p),
       H^b_{\cts}\bigl(P_n,\pi_t(E_n/p)\bigr)\right)
 \ \Longrightarrow\
 H^{a+b}_{\cts}\bigl(\mathbb G_n,\pi_t(E_n/p)\bigr).
\]
Its coefficients on the second page are finite, and a finite group has
finite cohomology in each degree on a finite module.  Each total-degree
diagonal has finitely many entries.  The abutment is therefore finite.
This reasoning also applies when $p$ divides the degree $n$ of the
constant-field extension; no averaging by the order of the Galois group
is used.
\end{proof}
